\documentclass[12pt]{amsart}
%   \overfullrule=5pt
\usepackage[active]{srcltx}
\usepackage{calc,amssymb,amsthm,amsmath,amscd, eucal,ulem}
\usepackage{alltt}
%\usepackage[left=1in,top=1in,right=1in,bottom=1in]{geometry}
\synctex=1
%\usepackage{mathtools}
\RequirePackage[dvipsnames,usenames]{color}
\let\mffont=\sf
\normalem
\input{mabliautoref.sty}
\input{xy}
\xyoption{all}
\input{kmacros3.sty}
\usepackage{tikz}
\usepackage{amsfonts, mathrsfs}
\usepackage{calligra}
\usepackage{stmaryrd} %power series bracket
%\usepackage{dcpic, pictexwd}
%\usepackage{pdfsync}
\usepackage[left=1.1in,top=1in,right=1.1in,bottom=1in]{geometry}
\usepackage{enumitem}

\numberwithin{equation}{theorem}

%\renewcommand{\thefootnote}{\fnsymbol{footnote}}
\newcommand{\D}{\displaystyle}
\newcommand{\til}{\widetilde}
\newcommand{\ol}{\overline}
\newcommand{\F}{\mathbb{F}}
\DeclareMathOperator{\E}{E}
\renewcommand{\:}{\colon}
\newcommand{\eg}{{\itshape e.g.} }
\renewcommand{\m}{\mathfrak{m}}
\renewcommand{\n}{\mathfrak{n}}
\newcommand{\Tt}{{\mathfrak{T}}}
\newcommand{\calC}{\mathcal{C}}
\newcommand{\RamiT}{\mathcal{R_{T}}}
\DeclareMathOperator{\edim}{edim}
\DeclareMathOperator{\length}{length}
\DeclareMathOperator{\monomials}{monomials}
\DeclareMathOperator{\Fitt}{Fitt}
\DeclareMathOperator{\depth}{depth}
\newcommand{\Frob}[2]{{#1}^{1/p^{#2}}}
\newcommand{\Frobp}[2]{{(#1)}^{1/p^{#2}}}
\newcommand{\FrobP}[2]{{\left(#1\right)}^{1/p^{#2}}}
\newcommand{\extends}{extends over $\Tt${}}
\newcommand{\extension}{extension over $\Tt${}}
\newcommand{\fg}{\pi_1^{\textnormal{\'{e}t}}} %Fundamental group
\DeclareMathOperator{\ns}{ns}
%\newcommand{\Spec}{\text{Spec}} %Spectrum
\newcommand{\pSpec}{\text{Spec}^{\circ}} %Puntured Spectrum
\newcommand{\etInCdOne}{\etale{} in codimension $1$}
\DeclareMathOperator{\DIV}{Div}
%\renewcommand{\char}{\charact}


\usepackage{setspace}
\usepackage{hyperref}
%\usepackage{enumerate}
\usepackage{graphicx}
\usepackage[all,cmtip]{xy}

\usepackage{verbatim}

\theoremstyle{theorem}
\newtheorem*{mainthm}{Main Theorem}
\newtheorem*{mainthma}{Main Theorem A}
\newtheorem*{mainthmb}{Main Theorem B}
\newtheorem*{mainthmc}{Main Theorem C}
\newtheorem*{atheorem}{Theorem}
%\newcommand{\myH}{h}
\newcommand{\sol}[1]{ \vskip 6pt{\color{blue} {\bf Solution:} #1} \vskip 6pt}

%The todo box!
\def\todo#1{\textcolor{Mahogany}%
{\footnotesize\newline{\color{Mahogany}\fbox{\parbox{\textwidth-15pt}{\textbf{todo:
} #1}}}\newline}}

%What is going on?  Why isn't \O a script O?!?
\renewcommand{\O}{\mathscr O}

\begin{document}
\title{In class worksheet \#2, getting a feel for Serre's condition ${\bf S}_n$}
\author{March 8th, 2017}
%\address{Department of Mathematics\\ University of Utah\\ Salt Lake City\\ UT 84112}
%\email{schwede@math.utah.edu}


%  \subjclass[2010]{14F18, 13A35, 14B05}

\maketitle

%\section{Reduction to characteristic $p > 0$}
%\chapter{Frobenius and Kunz's theorem}
%\bibliographystyle{skalpha}
%\bibliography{MainBib}

Our goal is to prove the following result.
\begin{proposition*}
Suppose that $(R, \fram)$ is a $d$-dimensional ${\bf S}_n$ ($n \leq d$) Noetherian equidimensional\footnote{This means for each minimal prime $\frq$ of $R$, $\dim R_{\frq} = \dim R$.} local ring with normalized dualizing complex $\omega_R^{\mydot}$.  Then $\dim \Supp \myH^{-i}(\omega_R^{\mydot}) \leq i - n$ for $i < d$.
\end{proposition*}

To prove this, you may use the following facts:
\begin{itemize}
\item{} Local duality, $\Hom_R(\myR \Hom_R(M, \omega_R^{\mydot}), E) \cong \myR \Gamma_{\fram}(M)$.
\item{} Since $\myH^{-i}(\omega_R^{\mydot})$ is a finitely generated $R$-module, $\Supp \myH^{-i}(\omega_R^{\mydot}) = V(\Ann_R(\myH^{-i}(\omega_R^{\mydot})))$
\item{} Since $R$ is equidimensional, $\omega_R = \myH^{-d}(\omega_R^{\mydot})$ has support equal to $\Spec R$.
\item{} Rings with dualizing complexes are catenary, and so every maximal chain of primes between two fixed primes has the same length.
\end{itemize}
\noindent
{\bf 1.}  For any prime $Q$ of $R$, show that $d = \dim R_Q + \dim R_Q$.  Give an example where this does not hold for a \emph{non-}equidimensional (but otherwise nice) ring.
\vskip 150pt
Now we prove the following weak version of the proposition.
\vskip 3pt
\noindent
{\bf 2.}  Suppose that $(R, \fram)$ is a $d$-dimensional Notherian ring with normalized dualizing complex $\omega_R^{\mydot}$.  Show that $h^{i}(\omega_R^{\mydot}) = 0$ for $i > 0$.
\newpage
\noindent
{\bf 3.}  Suppose that $Q \in \Spec R$ is a prime ideal such that $\dim R_Q = c$.  Find a normalized dualizing complex for $R_Q$.
\vskip 80pt
\noindent
{\bf 4.}  Suppose that $(R, \fram)$ is $\Serre_n$ ($n \leq d$) show that $\myH^{-i}(\omega_R^{\mydot}) = 0$ for $i < n$.
\vskip 125pt
\noindent
{\bf 5.}  Prove the proposition from the first page.
\vskip 3pt
\emph{Hint:} If $\dim \Supp \myH^{-i}(\omega_R^{\mydot}) > i - n$ for some $i < d$ choose an appropriate prime to localize at, then keep track of the \emph{normalized} dualizing complex.  

%\sol{Fix an $i$ to work with.  We need to show that $\dim \Supp \myH^{-i}(\omega_R^{\mydot}) \leq i - n$ for $i < d$.  Choose $Q$ a minimal prime of a component of the support of $\myH^{-i}(\omega_R^{\mydot})$ which has dimension $> i-n$.  Since $Q$ is a minimal prime of the support, $\dim R/Q > i-n$ and so $i - \dim R/Q < n$.  We know that $(\omega_R^{\mydot}[-\dim R/Q])_Q$ is a normalized dualizing complex for $R_Q$ which is also $\Serre_n$.  Thus
%\[
%\myH^{-j}((\omega_R^{\mydot}[-\dim R/Q])_Q) = 0 \text{ for $j < \min(n, \dim R_Q)$.}
%\]
%But by construction $R_Q$ is not Cohen-Macaulay (as otherwise its dualizing complex would live in a single degree but it does not), and so $\min(n, \dim R_Q) = n$.  Now, choosing $j = i - \dim R/Q < n$, we see that
%\[
%0 = \myH^{-j}((\omega_R^{\mydot}[-\dim R/Q])_Q) = \myH^{-i + \dim R/Q - \dim R/Q}((\omega_R^{\mydot})_Q) = \myH^{-i}((\omega_R^{\mydot})_Q) \neq 0
%\]
%a contradiction.
%}
\newpage
Let's do some other applications of dualizing complexes.  Recall that a local ring $(R, \fram)$ is called $F$-injective if for every $i \geq 0$, the induced Frobenius action
\[
H^i_{\fram}(R) \to H^i_{\fram}(F_* R)
\]
is injective.  A Noetherian ring in general is called $F$-injective if for each maximal ${\fram} \in \Spec R$, $R_{\fram}$ is $F$-injective.
\vskip 3pt
\noindent
{\bf 6.}  Suppose that $(R, \fram)$ is an $F$-finite Noetherian local ring.  Show that for each $Q \in \Spec R$, $R_Q$ is $F$-injective.
\vskip 200pt
\noindent
{\bf 7.}  Suppose that $R$ is an $F$-finite Noetherian ring with dualizing complex $\omega_R^{\mydot}$.  Show that $R$ is $F$-injective if and only if the induced map $F_* \omega_{R^{\mydot}} \cong \omega_{F_* R}^{\mydot} \to \omega_R^{\mydot}$ surjects on all cohomology (\ie, $F_* \myH^i (\omega_R^{\mydot}) \to \myH^i \omega_R^{\mydot}$ for all $i$).


\end{document}


